% =====================================================================
% Rozdział 4 – Implementacja systemu (~16 str.)
% Kolejność zgodna z przepływem danych: model -> kalibracja i LUT ->
% węzeł -> UDP -> fuzja -> ROS 2. Narzędzia (ESP-IDF, Edge Impulse,
% TFLite Micro, ESP-NN, Python, OpenCV, ROS 2) opisywać tam, gdzie są
% użyte, z wersjami. Problemy techniczne krótko w miejscu wystąpienia
% (problem -> przyczyna -> rozwiązanie).
% =====================================================================
\chapter{Implementacja systemu}
\label{ch:implementacja}

% ---------------------------------------------------------------------
\section{Zbiór danych i~trening modelu}
\label{sec:zbior-trening}
% - Narzędzia akwizycji:
%   - firmware ESP32-CAM + collect_dataset.py przez UART (Python wysyła
%     „c”, ESP odpowiada „IMG:<base64 JPEG>”), sterowanie klawiaturą,
%     nazewnictwo plików,
%   - problemy: reset płytki przez DTR/RTS (ESP32-CAM-MB), zły model
%     płytki w kodzie (OV3660),
%   - drugie narzędzie: projekt ESP-IDF na ESP32-S3 – podgląd po Wi-Fi
%     + przycisk capture.
% - Dlaczego zdjęcia 320x240, a nie 96x96: zachowanie oryginałów,
%   ponowne przetworzenie (LUT), wygodniejsze etykietowanie.
% - Zbiór: butelka + tło; liczebność (plan ok. 150 + ok. 65) – faktyczne
%   liczby ⏳; różnorodność (odległości, położenie w kadrze także przy
%   krawędziach, oświetlenie, tła, typy butelek).
% - Przetwarzanie: 320x240 -> LUT (\ref{sec:lut}) -> 96x96 -> PNG
%   (bezstratnie). Model trenowany na obrazach po LUT – widzi dokładnie
%   to, co węzeł.
% - Etykietowanie w Edge Impulse, podział trening/test; dane publiczne
%   tylko jako uzupełnienie; walidacja potoku na projekcie „Beer vs. Can”.
% - Trening: wejście 96x96x1, FOMO na MobileNetV2 alpha 0.1, epoki,
%   learning rate, augmentacja ⏳; kwantyzacja int8; eksport jako
%   biblioteka C++ bez EON (dlaczego – \ref{sec:oprogramowanie-wezla}).
% - Rys.: przykładowe zdjęcia przed i po LUT.
% - Tab.: liczebność zbioru i parametry treningu.
% Stan: ✅ / ⏳ (liczby, parametry)


% ---------------------------------------------------------------------
\section{Kalibracja parametrów wewnętrznych kamer}
\label{sec:kalibracja-wewn}
% - Procedura: szachownica (np. 9x6 narożników) na sztywnym podłożu,
%   30–40 zdjęć na kamerę w trybie pracy (QVGA), wzorzec w całym kadrze,
%   zwłaszcza w rogach.
% - Ostrość obiektywu ustawiona raz i zablokowana przed kalibracją.
% - Python + OpenCV: findChessboardCorners, cornerSubPix, calibrateCamera
%   (CALIB_FIX_K3 przy umiarkowanej dystorsji) – wersje bibliotek.
% - Wyniki dla każdej kamery: K, współczynniki dystorsji, RMS reprojekcji
%   (cel < 0,5 px), wyliczony HFOV.
% - Tab.: wyniki kalibracji 6 kamer.
% - Rys.: wykryte narożniki na zdjęciu.
% Stan: ⏳ (kalibracja wszystkich kamer)


% ---------------------------------------------------------------------
\section{Korekcja dystorsji z~użyciem tablicy LUT}
\label{sec:lut}
% - Dwa podejścia: metoda A – sieć na obrazie zniekształconym, poprawiany
%   tylko centroid na PC; metoda B2 – obraz prostowany na węźle przed
%   siecią.
% - Dlaczego B2: model widzi obraz bez dystorsji (spójny między kamerami),
%   obiekty przy krawędziach nie są zniekształcone, koszt na węźle mały.
% - Generowanie LUT (Python, OpenCV: getOptimalNewCameraMatrix,
%   initUndistortRectifyMap): jedna tablica robi naraz korekcję dystorsji
%   i skalowanie 320x240 -> 96x96 (squash 4:3 -> 1:1 przez osobne
%   skalowanie fx, fy).
% - Wspólne K_new dla wszystkich kamer – jeden model działa na każdej;
%   macierz K96 (z poprawką na środek piksela) używana na PC.
% - Format: uint16, Q9.7, 0xFFFF = brak piksela; 2 tablice po 9216
%   wpisów ≈ 36 kB flash na kamerę (lut_camN.h).
% - Implementacja na węźle: interpolacja dwuliniowa w arytmetyce
%   całkowitej, pola spoza obrazu wypełniane stałą; koszt 4,5 ms
%   (ok. 4% pętli).
% - Test zgodności: ta sama klatka w Pythonie (apply_lut.py) i na ESP –
%   wynik identyczny lub ±1.
% - Pułapki: osobna LUT na kamerę, zmiana ostrości/trybu sensora
%   unieważnia LUT, ten sam preprocessing w treningu i firmware.
% - Rys.: obraz przed i po LUT (proste linie).
% - Rys.: schemat generowania LUT.
% - Listing: funkcja LUT (fragment).
% Stan: ✅ (metoda), ⏳ (LUT dla wszystkich kamer)


% ---------------------------------------------------------------------
\section{Oprogramowanie węzła}
\label{sec:oprogramowanie-wezla}
% - Środowisko: ESP-IDF 5.3.1 (FreeRTOS, komponenty, menuconfig/sdkconfig),
%   VS Code + rozszerzenie ESP-IDF, idf.py. Dlaczego ESP-IDF zamiast
%   Arduino: pełna kontrola konfiguracji, dostęp do ESP-NN; przebieg
%   migracji.
% - Struktura projektu: main/, komponenty (kamera, biblioteka Edge
%   Impulse), pliki lut_camN.h.
% - Kluczowe ustawienia sdkconfig: PSRAM octal 80 MHz, cache danych
%   64 kB / linia 64 B, tick 1000 Hz, wyłączona kontrola sterty, partycja
%   „Single factory app (large)”.
% - Kamera: komponent esp32-camera, QVGA, grayscale, XCLK 20 MHz,
%   2 bufory, GRAB_LATEST.
% - Pętla główna: klatka -> LUT -> przygotowanie wejścia -> inferencja
%   -> centroidy -> UDP.
% - Inferencja: Edge Impulse SDK -> TensorFlow Lite Micro (interpreter,
%   arena) -> jądra ESP-NN (SIMD). Arena w PSRAM (po starcie Wi-Fi
%   zostaje ok. 66 kB RAM wewnętrznego).
% - Trwały interpreter: łatka SDK Edge Impulse – interpreter i arena
%   tworzone raz zamiast co klatkę (13,7 -> 0,8 ms); powtarzać po każdym
%   eksporcie.
% - Wyniki FOMO: centroidy w układzie 96x96, próg pewności.
% - Problemy (problem -> przyczyna -> rozwiązanie):
%   - model EON: crash przy pierwszej inferencji -> eksport bez EON,
%   - ICE GCC 13.2 w nagłówkach Edge Impulse -> main.cpp na -O0,
%     gorące pętle z __attribute__((optimize("O2"))),
%   - kontrola sterty ok. 400 ms/klatkę -> wyłączona,
%   - zapełniona partycja 1 MB -> większa partycja,
%   - zmiana formatu kamery wymaga zmiany ei_camera_capture().
% - Rys.: diagram pętli węzła.
% - Tab.: kluczowe ustawienia konfiguracji.
% Stan: ✅


% ---------------------------------------------------------------------
\section{Komunikacja UDP i~znakowanie czasem}
\label{sec:udp}
% - Zawartość pakietu: identyfikator węzła, klasa, centroid (u, v
%   w 96x96), pewność, czas wykonania zdjęcia.
% - Numer sekwencyjny klatki – liczenie utraconych pakietów
%   (\ref{sec:wydajnosc}). ❓ czy jest/dodać.
% - Format binarny vs tekstowy – co i dlaczego. ⏳ faktyczny format.
% - Znacznik czasu: wcześniej moment odebrania klatki, teraz czas zdjęcia
%   ze sterownika kamery – różnica rzędu wieku klatki.
% - Zegar ESP liczy od startu -> SNTP na węzłach albo estymacja
%   przesunięcia zegara na PC. ❓ wybór metody.
% - Obsługa Wi-Fi: łączenie, ponowne łączenie po utracie sieci.
% - Odbiór na komputerze (Python, fomo_udp_receiver.py lub węzeł ROS 2):
%   gniazdo, rozpoznanie węzła, dekodowanie, przeliczenie czasu węzła na
%   czas PC, bufor ostatnich detekcji z każdej kamery w oknie czasowym,
%   wykrywanie utraty pakietów i nieaktywnych węzłów (timeout).
% - Tab.: pola pakietu (nazwa, typ, rozmiar, opis).
% Stan: ✅ / ⏳ / ❓


% ---------------------------------------------------------------------
\section{Kalibracja zewnętrzna i~fuzja detekcji}
\label{sec:kalibracja-zewn-fuzja}
% - Orientacja i położenie kamer względem układu robota; ponieważ
%   wyznaczane są tylko kierunki, najważniejszy jest kąt obrotu;
%   przesunięcie ma znaczenie dla bliskich obiektów (paralaksa).
% - Metoda ❓: z geometrii obudowy (model OpenSCAD) + korekta / obiekt
%   referencyjny w znanych kierunkach / obszary nakładania.
% - Przeliczenie detekcji: centroid (u, v) -> x_n = (u - cx96)/fx96
%   (obraz już prosty) -> alfa = atan(x_n) -> theta = alfa + psi_kamery.
%   Dokładność wynikająca z siatki FOMO (komórka 8 px).
% - Algorytm fuzji:
%   - wejście: detekcje (azymut, klasa, pewność, czas) ze wszystkich kamer,
%   - dopasowanie w czasie: okno czasowe; kompensacja obrotu robota
%     z odometrii,
%   - asocjacja: detekcje z sąsiednich kamer w obszarze nakładania
%     o różnicy kąta poniżej progu -> ten sam obiekt,
%   - łączenie: średnia ważona pewnością lub kamera, w której obiekt jest
%     bliżej środka kadru,
%   - wyjście: lista obiektów (azymut, klasa, pewność, czas),
%   - parametry (próg kąta, okno czasu, próg pewności) i ich dobór;
%     złożoność pomijalna.
% - Rys.: stanowisko kalibracyjne.
% - Rys.: obiekt w obszarze nakładania widziany przez 2 kamery ->
%   1 obiekt po fuzji.
% - Pseudokod algorytmu.
% Stan: ❓ (metoda kalibracji) / ⏳ (fuzja)


% ---------------------------------------------------------------------
\section{Integracja z~ROS~2}
\label{sec:ros2}
% - ROS 2 krótko w kontekście użycia: węzły, tematy, publish/subscribe.
%   Wersja i system. ⏳
% - Węzły ESP NIE używają micro-ROS – zwykłe UDP, do ROS 2 publikuje
%   komputer (prostota, mniej RAM na węźle, brak agenta).
% - Węzeł ROS 2 na komputerze: odbiór -> fuzja -> publikacja.
% - Tematy: surowe detekcje z każdej kamery i wynik fuzji. ❓ nazwy.
% - Typ komunikatu: własny .msg (azymut, klasa, pewność, czas) lub
%   vision_msgs. ❓
% - Znacznik czasu w nagłówku = czas zdjęcia, nie czas publikacji.
% - Wizualizacja w RViz: strzałki/markery kierunków wokół robota.
% - Rys.: graf węzłów i tematów (rqt_graph).
% - Rys.: zrzut z RViz.
% Stan: ⏳

